\begin{ouverture}[mesure]{1}{Mesurer et raisonner}{Outils de la démarche scientifique}{Observer · Modéliser · Expérimenter · Expliquer}
\objectif{Choisir les unités et donner un ordre de grandeur.}
\objectif{Interpréter la dispersion de mesures indépendantes.}
\objectif{Présenter un résultat et discuter sa compatibilité avec une référence.}
\prerequis{Proportionnalité, conversions et lecture de graphiques.}
\end{ouverture}
\essentiel
\essSection{1}{Une grandeur, un nombre et une unité}
Une mesure associe une valeur à une grandeur : longueur, durée, masse, température. Écrire « 12 » ne suffit pas à décrire une longueur. Le symbole de l'unité ne prend pas de pluriel. Avant un calcul, choisir un ensemble cohérent d'unités.
\begin{definition}[Écriture scientifique]Un nombre positif s'écrit $a\times10^p$ avec $1\leq a<10$ et $p$ entier. Les préfixes milli, micro et nano correspondent à $10^{-3}$, $10^{-6}$ et $10^{-9}$.\end{definition}
Un ordre de grandeur sert à repérer un résultat invraisemblable. Une règle graduée au millimètre ne permet pas, à elle seule, de distinguer des longueurs au micromètre.
\begin{exemple}[Convertir avant de diviser]$V=\qty{25,0}{mL}=\qty{0,0250}{L}$ ; $m=\qty{0,450}{g}$. Le quotient $m/V$ vaut $\qty{18,0}{g.L^{-1}}$.\end{exemple}
\essSection{2}{Mesures répétées et dispersion}
Même avec le même protocole, des mesures indépendantes diffèrent légèrement : lecture, remplissage, environnement. Refaire le même calcul sur la même mesure ne constitue pas une nouvelle mesure indépendante.
\begin{definition}[Moyenne et écart-type]La moyenne $\bar x$ situe la série. L'écart-type expérimental $s$ décrit sa dispersion. Un histogramme rassemble les valeurs par intervalles ; sa forme dépend aussi du choix de ces intervalles.\end{definition}
Pour $N$ mesures indépendantes réalisées dans les mêmes conditions, on peut estimer l'incertitude-type sur la moyenne par $u(\bar x)=s/\sqrt N$. L'écart-type $s$ concerne la dispersion des valeurs individuelles ; il ne faut pas le confondre avec $u(\bar x)$.
\schema[corrige]{histogramme}
\essSection{3}{Écrire et discuter un résultat}
On conserve des chiffres pendant le calcul, puis on arrondit le résultat final. Lorsqu'une incertitude est donnée, la valeur et l'incertitude s'écrivent au même rang décimal.
\begin{exemple}[Une conclusion mesurée]Pour $\bar V=\qty{10,03}{mL}$ et $u(\bar V)=\qty{0,02}{mL}$, écrire ces deux valeurs avec leur unité. Une référence de \qty{10,00}{mL} est proche à l'échelle de quelques incertitudes-types ; une référence de \qty{10,50}{mL} est nettement éloignée. En Seconde, on discute qualitativement cette comparaison.\end{exemple}
\begin{attention}[Ce qu'une répétition ne corrige pas]Une balance mal réglée peut décaler toute une série. Répéter réduit l'effet des fluctuations sur la moyenne, mais ne supprime pas ce décalage. Une incertitude-type n'est pas une borne absolue de l'erreur.\end{attention}
\essSection{4}{Expérimenter et modéliser}
\begin{experience}[Comparer deux verreries]Peser l'eau délivrée par une éprouvette puis par une pipette, après tare du récipient. Recommencer dix fois en vidant et remplissant à nouveau. À température donnée, convertir les masses en volumes avec la masse volumique fournie. Comparer les moyennes et les dispersions, puis discuter le choix de la verrerie.\end{experience}
\begin{savoirfaire}[Tableur]Saisir une mesure par cellule. Calculer la moyenne et l'écart-type d'échantillon, choisir des classes contiguës et de même largeur, puis produire l'histogramme. L'axe horizontal porte la grandeur et l'unité ; l'axe vertical porte l'effectif. Ne pas relier les barres comme une trajectoire.\end{savoirfaire}
\exercices
\rubrique{Prendre confiance}
\begin{blocLivre}{ex-1-1}\begin{exercice}[1]{Changer de préfixe}
Un capteur détecte un déplacement de \qty{450}{\micro m} et une durée de \qty{2,5}{ms}.
\begin{questions}
\item Exprimer ces deux valeurs dans les unités SI.
\item Écrire \qty{0,000072}{m} en écriture scientifique.
\end{questions}
\end{exercice}\end{blocLivre}
\begin{blocLivre}{ex-1-2}\begin{exercice}[1]{Un résultat sans unité}
Une élève écrit : « Le volume vaut 25 et la masse 19,8 g ».
\begin{questions}
\item Pourquoi le premier résultat est-il incomplet ?
\item Compléter si le volume a été lu en millilitres et convertir en litres.
\end{questions}
\end{exercice}\end{blocLivre}
\begin{blocLivre}{ex-1-3}\begin{exercice}[1]{Lire un instrument}
Une règle porte des traits espacés de \qty{1}{mm}. On mesure une longueur proche de \qty{8,4}{cm}.
\begin{questions}
\item Choisir entre 8,4 cm et 8,40000 cm pour rendre compte de cette lecture.
\item Citer deux causes possibles de variation entre lectures.
\end{questions}
\end{exercice}\end{blocLivre}
\begin{blocLivre}{ex-1-4}\begin{exercice}[1]{Une moyenne simple}
Quatre durées indépendantes valent 2,1 ; 2,0 ; 2,2 ; 2,1 s.
\begin{questions}
\item Calculer la moyenne.
\item Pourquoi ne faut-il pas supprimer 2,2 s uniquement parce que cette valeur est la plus grande ?
\end{questions}
\end{exercice}\end{blocLivre}
\rubrique{S’entraîner}
\begin{blocLivre}{ex-1-5}\begin{exercice}[2]{Deux distributions}
La figure présente deux séries de mesures d’un volume nominal de \qty{10,0}{mL}.
\schema{histogramme}
\begin{questions}
\item Quelle série est la moins dispersée ?
\item Une faible dispersion garantit-elle l’absence de décalage systématique ?
\end{questions}
\end{exercice}\end{blocLivre}
\begin{blocLivre}{ex-1-6}\begin{exercice}[2]{Estimer une incertitude}
Pour 16 mesures indépendantes, on obtient $\bar V=\qty{24,96}{mL}$ et $s=\qty{0,12}{mL}$.
\begin{questions}
\item Calculer $u(\bar V)=s/\sqrt N$.
\item Présenter la moyenne et son incertitude-type.
\end{questions}
\end{exercice}\end{blocLivre}
\begin{blocLivre}{ex-1-7}\begin{exercice}[2]{Arrondir à la fin}
On calcule $v=0,625/0,48$ en unités SI.
\begin{questions}
\item Calculer sans arrondir les étapes intermédiaires.
\item Proposer une écriture à deux chiffres significatifs.
\end{questions}
\end{exercice}\end{blocLivre}
\begin{blocLivre}{ex-1-8}\begin{exercice}[2]{Comparer à une référence}
Un groupe obtient $\bar \rho=\qty{0,79}{g.mL^{-1}}$, avec $u=\qty{0,01}{g.mL^{-1}}$. Deux références sont 0,789 et \qty{1,00}{g.mL^{-1}}.
\begin{questions}
\item Laquelle paraît compatible à l’échelle de l’incertitude ?
\item Pourquoi cela ne prouve-t-il pas l’identité du liquide ?
\end{questions}
\end{exercice}\end{blocLivre}
\rubrique{Raisonner et approfondir}
\begin{blocLivre}{ex-1-9}\begin{exercice}[3]{Un quotient trompeur}
Un calcul donne $m/V=250/0,50=500$ à partir de $m=\qty{250}{mg}$ et $V=\qty{0,50}{L}$.
\begin{questions}
\item Associer l’unité correcte à 500.
\item Exprimer le résultat en \unit{g.L^{-1}}.
\end{questions}
\end{exercice}\end{blocLivre}
\begin{blocLivre}{ex-1-10}\begin{exercice}[3]{Répéter utilement}
Deux élèves disposent de la même mesure de masse. Le premier la saisit dix fois ; le second effectue dix nouvelles pesées après avoir vidé puis rempli le récipient.
\begin{questions}
\item Quelle démarche permet d’étudier la variabilité du remplissage ?
\item Le second doit-il conserver la tare et la température sous contrôle ?
\end{questions}
\end{exercice}\end{blocLivre}
\begin{blocLivre}{ex-1-11}\begin{exercice}[3]{Améliorer une méthode}
On chronomètre à la main une seule oscillation de durée proche de 0,8 s. Le temps de réaction gêne la mesure.
\begin{questions}
\item Proposer une méthode avec le même chronomètre.
\item Dix oscillations durent 8,16 s ; déterminer la période moyenne.
\end{questions}
\end{exercice}\end{blocLivre}
\begin{blocLivre}{ex-1-12}\begin{exercice}[3]{Histogramme et classes}
On mesure 9,8 ; 9,9 ; 10,0 ; 10,0 ; 10,1 ; 10,2 mL. Les classes sont $[9,75;9,95[$, $[9,95;10,15[$ et $[10,15;10,35]$.
\begin{questions}
\item Donner les trois effectifs.
\item Calculer la moyenne et expliquer ce que l’histogramme apporte de plus.
\end{questions}
\end{exercice}\end{blocLivre}

